\documentclass[answers]{exam}
\usepackage{../../mypackages}
\usepackage{../../macros}

% Pour le tableau
\usepackage{array}
\usepackage{longtable}
\usepackage{multirow}
\usepackage[table]{xcolor}

\title{Barème de notation des cahiers}
\author{N. Bancel}
\date{2 Décembre 2025}

% Mise en forme des solutions en bleu
\SolutionEmphasis{\color{blue}}
\renewcommand{\solutiontitle}{\noindent}

\definecolor{lavande}{RGB}{180,190,255}

% Colonne p{} alignée à gauche et qui wrappe
\newcolumntype{L}[1]{>{\raggedright\arraybackslash}p{#1}}

\begin{document}

\textbf{Collège Lycée Suger}
\hfill
\textbf{Mathématiques} \\

\textbf{Année 2025-2026}
\hfill
\textbf{Classe de 6ème} \par

{\let\newpage\relax\maketitle}


\renewcommand{\arraystretch}{1.7}


\centering
\begin{longtable}{|L{0.15\textwidth}|L{0.3\textwidth}|L{0.5\textwidth}|c|}
\hline
\rowcolor{lavande}
\textbf{Catégorie} & \textbf{Critère évalué} & \textbf{Description} & \textbf{Points} \\[2pt]
\hline

% ---------------- Cahier d'exercices ----------------
\multirow{3}{*}{%
  \parbox[t]{0.15\textwidth}{\raggedright Cahier d’exercices}%
}
  & Titres des exercices présents et soignés (+ utilisation de couleur)
  & Les exercices comportent un titre clair, souligné ou encadré. Et une référence claire à l'endroit où retrouver l'exercice (« Exercice n°34 page 127 – Le livre scolaire », « Exercice d'application n°3 – Lutin », « Exercice n°2 – Fiche d'exercice »).
  & 2 \\
\cline{2-4}

  & Ordres de grandeurs notés
  & Les ordres de grandeurs donnés en début d'heure ont bien été notés.
  & 1 \\
\cline{2-4}

  & Corrections notées (ou rattrapées en cas d'absence)
  & Toutes les corrections sont recopiées, complétées ou rattrapées en cas d’absence, de manière lisible. Même si vous avez juste, j'attends de vous que vous preniez la correction. Pour progresser sur la rédaction.
  & 3 \\[6pt]
\hline  

% ---------------- Lutin ----------------
\multirow{5}{*}{Lutin}
  & Feuilles rangées dans l’ordre
  & Les chapitres sont dans l'ordre. Et les feuilles sont classées dans l’ordre attendu : Recto : Page 1. Verso : Page 2. Puis nouvelle pochette transparente : Page 3 au recto, Page 4 au verso etc.
  & 1 \\
\cline{2-4}

  & 1-2 feuilles max par pochette
  & Chaque feuille imprimée dans un feuillet transparent (1-2 feuilles par pochette transparente). Je ne veux pas voir toutes les feuilles d'un chapitre dans une seule et même pochette transparente.
  & 1 \\
\cline{2-4}

  & Interrogations + calcul mental faciles à retrouver
  & Les évaluations de calcul mental sont regroupées et rapidement retrouvables.
  & 2 \\
\cline{2-4}

  & Cours complet (rien de manquant)
  & Aucune feuille manquante : l’ensemble du cours est présent dans le lutin.
  & 2 \\
\cline{2-4}

  & Leçons, phrases, définitions recopiées correctement
  & Le travail est propre : tracés à la règle ou compas, titres alignés, soin général. La prise de notes est complète. Il ne manque pas de définition, de schéma.
  & 2 \\[6pt]
\hline

% ---------------- Général ----------------
\multirow{5}{*}{Général}
  & Organisation générale (titres soulignés, couleurs…)
  & Le cahier et le lutin sont organisés : titres visibles, couleur, structure claire.
  & 2 \\
\cline{2-4}

  & Écriture soignée et lisible
  & L’écriture est propre, lisible, régulière.
  & 1 \\
\cline{2-4}

  & Peu ou pas de ratures
  & Les ratures sont limitées et propres, sans nuire à la lecture.
  & 1 \\
\cline{2-4}

  & Pas de feuilles volantes
  & Aucune feuille glissée en vrac ; tout est rangé ou collé.
  & 1 \\
\cline{2-4}

  & Présence de prénom / nom // Identification facile du lutin
  & On sait très clairement à qui appartient le lutin, et à qui appartient le cahier d'exercice. Il y a une page de garde qui indique clairement qu'il s'agit de votre matériel de Mathématiques
  & 1 \\[6pt]
\hline

% ---------------- Total ----------------
\rowcolor{green!20}
 & \textbf{Total} & \textbf{Note maximale possible} & \textbf{20} \\
\hline

\end{longtable}




\section{Conseils}

\raggedright
\noindent\textbf{\large Gestion du cahier d'exercice}\par\vspace{0.5\baselineskip}

\begin{compactitem}
  \item A chaque début de chapitre, écrivez en gros, en rouge le titre du chapitre. Encadrez-le. Tout ce qui sera écrit seront des exercices en lien avec le Chapitre
  \item Indiquez clairement à quel exercice une correction est associée. On écrit, dans une couleur bien identifiable (rouge par exemple, souligné, ou encadré) : "Exercice N°26 Page 127 du Livre Scolaire". Il faut que vous puissiez retrouver facilement l'exercice auquel vous faites référence.
  \item Vous pouvez même coller l'énoncé dans le cahier. 
  \item Même quand vous avez juste : si vous voyez que ma rédaction est plus complète, j'attends que vous la notiez
\end{compactitem}

\vspace{1em}
\noindent\textbf{\large Gestion du lutin}\par\vspace{0.5\baselineskip}

\begin{compactitem}
  \item Je ne veux pas voir des Chapitres dans le désordre, où le chapitre 7 arrive avant le chapitre 2.
  \item Respectez la numérotation des pages : Page 1 au recto, Page 2 au verso. Puis nouvelle pochette transparente : Page 3 au recto, Page 4 au verso etc.
  \item L'imprimante a eu quelques problème, en imprimant dans l'autre sens des versos. Je vous conseille de ré imprimer les feuilles qui posent problème. Ou demandez-moi. 
  \item Je vous conseille de coller vos interrogations de calcul mental sur des feuilles simples ou doubles. Et stockez ces feuilles à la fin de votre lutin. 
  \item Il est attendu que vos Devoirs sur table soient aussi stockés dans votre lutin. 
  \item Prise de note : 
  \begin{compactitem}
    \item Les figures doivent être faites au crayon à papier.
    \item Le reste de la prise de note se fait au stylo.
    \item Des légendes doivent figurer sur les figures.
    \item On ne fait pas de tracé à main levée, on réalise ses figures avec compas, rapporteur, règle, équerre
  \end{compactitem}
\end{compactitem}

\vspace{1em}
\noindent\textbf{\large Général}\par\vspace{0.5\baselineskip}

\begin{compactitem}
  \item Evitez les ratures - vos cahiers sont parfois illisibles
  \item Soignez votre écriture
  \item Il ne faut pas qu'il y ait de feuille volante. Tout doit être soit collé, soit rangé dans le lutin. 
\end{compactitem}


\end{document}
